The normal vector to the plane is $(1,1,1)$. A perpendicular line through $P$ therefore has points $(t,t,t)$. Intersecting with the plane gives $3t=1$, so the foot is $H=(\dfrac{1}{3},\dfrac{1}{3},\dfrac{1}{3})$. It is the average of the triangle\textquotesingle s three vertices and hence lies inside that triangle. The perpendicular distance is $|PH|=\sqrt{3\cdot\dfrac{1}{9}}=\boxed{\dfrac{1}{\sqrt3}}$. This identifies the foot by perpendicularity and intersection, without a visual guess.
\begin{asy}
import graph;
usepackage("amsmath");
import geometry;
size(7cm);
pair A=(3,0),B=(-1.5,1),C=(0,3),P=(0,0),H=(.5,4/3);draw(A--B--C--cycle);draw(P--A);draw(P--B);draw(P--C);draw(P--H,dashed);label("$A$",A,(1,0));label("$B$",B,(-1,0));label("$C$",C,(0,1));label("$P$",P,(0,-1));label("$H$",H,(1,1));
\end{asy}

